% Jeffrey's abstract, introduction and methods (2026-08-30 paste) with the FIXLIST_v12 corrections applied.
% Every change is listed at the end of this file under CHANGELOG. Sentences not listed there are verbatim.
\begin{abstract}

Sparse Mixture-of-Experts models contain a fixed pool of experts at each MoE layer. During both prompt prefill and autoregressive generation, a learned router selects a small subset of experts for each token, based on its current hidden representation.

Prior work correctly finds that prompt-side routing can converge on one shared set of experts across domains. We show that generation gives a different answer.

We analyze five checkpoints across the Qwen3.5 and GLM-4.7-Flash families, prefill on all five and generation on the two Qwen3.5-35B-A3B checkpoints. In the HauhauCS fine-tune of Qwen3.5-35B-A3B-Instruct, we test whether the selected expert sets follow the local subject of an answer by packing multiple subjects into one generation, letting the model regroup them across prompts, controlling for position, and comparing them with the same questions asked independently. We repeat the subject test on the official instruction-tuned checkpoint from which the HauhauCS model was derived.

Position-resolved prefill routing primarily follows the wording of the question, while pooling across the prompt produces the shared usage pattern reported in prior work. Across packed and independently asked versions of identical questions, generation-based expert sets identify the subject at 95 to 96\% in Qwen3.5-35B when the answer's register is held constant (0.70 over all design pairs, 0.54 when the register changes), compared with 15 to 17\% for prefill and 6\% by chance. Generation-time expert sets remain stable within a subject, change sharply at subject boundaries, and recur for the same subject across positions, prompt design, and the official and fine-tuned checkpoints. On the official instruction-tuned checkpoint, generation-time sets pair prompts by subject in 11 of 12 cases where prefill pairs 4 of 12. For short mixed prompts, prefill carries the wording of the question, while generation carries the subject of the answer.

\end{abstract}

\section{Introduction}

Expert specialization in sparse Mixture-of-Experts models is commonly inferred from which experts are selected for an input. That inference requires selective recurrence. A shared expert that appears across every subject may be important, but it is not specialized by subject. If the same expert set returns for the same subject after its wording, position, and surrounding context change, while different subjects recruit different sets, then the router is organizing computation by local subject. We use this subject-conditioned recurrence as the operational signature of specialization.

Prior work presents two views of expert organization. At the last prompt token of MMLU questions, the Standing Committee analysis finds that a small, fixed coalition carries most of the routing mass across nearly all domains \citep{wang2026standingcommittee}. Herbst et al. find that individual experts probe as fine-grained task specialists, such as experts associated with closing brackets or function definitions, rather than broad domain modules \citep{herbst2026expertstrikesback}. Ye et al. instead find that individual experts are polysemantic while their paths across layers become monosemantic \citep{ye2026polysemantic}. Do et al. identify domain-associated experts from question-token routing and show that upweighting them can improve domain accuracy across ten models \citep{do2026domainexperts}. These findings place specialization at different units: individual experts, collaborative sets, or cross-layer paths.

The distinction between prompt prefill and text generation creates another divide. Wang, Hayou, and Nalisnick find that semantically different prompts can converge on nearly identical expert usage during prefill and diverge during generation. Prompt-level routing therefore does not predict rollout-level routing, and complete trajectories are needed to characterize specialization \citep{wang2026myth}. Their result establishes the divergence, but leaves its organization unknown. The generated routes could follow the wording of the original question, move with token position, drift gradually through the answer, or separate according to the answer's local subject. We ask which of these possibilities explains the rollout.

We test this question across five checkpoints from the Qwen3.5 and GLM-4.7-Flash families. The primary controlled experiment uses the HauhauCS fine-tune of Qwen3.5-35B-A3B-Instruct. We pack questions from multiple subjects into a single prompt, generate one answer containing each subject, use the model's own regrouping of that answer to place subjects at different positions across prompts, and compare each generated segment with the same question asked independently. The same identification procedure is then applied to the prefill rows of those questions. This places the same semantic material on both sides of the prefill to generation boundary.

The design separates subject from four competing explanations. The model's regrouping places the same subject at different token positions and different subjects at matched positions. Restricting the candidate set to the preceding, current, and following segments tests gradual drift. Comparing packed and independently asked versions changes prompt design and surrounding context. Holding the voice of the packed answer constant separates subject from register, while a second design allows register to vary. We repeat a smaller subject-matching experiment on the official instruction-tuned checkpoint to test whether the organization was already present before the HauhauCS fine-tune. The remaining checkpoints provide the prefill replication and model-family scope.

We reproduce the shared prefill pattern reported in prior work. When prefill routing is pooled over each question, one expert wins 18 of 20 subjects in Qwen3.5-35B. Position-resolved prefill routing primarily follows the wording and local form of the question: across packed and independently asked versions of the same questions, its expert sets identify the subject only 15 to 17\% of the time. Generation gives a different answer. Expert-set similarity identifies the subject across those prompt designs at 95 to 96\%, against 6\% under shuffled labels.

The generation signal also survives the positional and temporal controls. The same subject at different positions matches more strongly than different subjects at the same position in nine of ten comparisons. When the candidates are restricted to the previous, same, and next subjects, expert-set similarity still selects the correct subject 83\% of the time, against 35\% by chance. The sets therefore change at subject boundaries rather than sliding gradually through the answer. On the official instruction-tuned checkpoint, generation pairs prompts by subject in 11 of 12 cases, while prefill does so in 4 of 12. Recurring legal and medical experts also survive changes in dialect and the official to fine-tuned checkpoint boundary.

These results place the apparently conflicting findings of prior work into two routing regimes. For short mixed prompts, pooled prefill routing exposes a shared usage pattern, while position-resolved prefill routing remains closely tied to question form. During generation, local subjects recruit reproducible expert sets that remain stable within a subject, turn over at its boundary, and return when the subject reappears under another prompt design. This distinction is scoped rather than universal: on a long, coherent technical prompt, prefill and generation share much of the same expert set. The separation is strongest when a short question initiates a substantially longer answer organized around its subject.

We make four contributions. First, we replicate the shared prefill result on five checkpoints across two model families and two gating mechanisms, three of which no prior study had tested. Second, we define subject specialization through cross-design recurrence and provide a model-agnostic identification protocol with wording, position, drift, and shuffled-label controls. Third, we show that generation-time expert sets track the local subject of an answer at 95 to 96\% identification accuracy, while prefill rows of the same questions reproduce the shared-routing result. Fourth, we show that this organization survives changes in position, prompt design, and the official instruction-tuned to fine-tuned checkpoint boundary. Register is a second organizer: subject survives when register is held constant, and a change of register moves the set as much as a change of subject does, lowering identification from 0.85 to 0.54. Expert specialization is visible when it is measured at the level of recurring generation-time sets.

\section{Methods}

\paragraph{Models and captures.} We study five checkpoints. The primary architecture is Qwen3.5-35B-A3B, with 256 routed experts per layer, top-8 routing, and 40 MoE layers, in two versions: the official instruction-tuned checkpoint (Q8\_0, ggml-org) and the HauhauCS uncensored fine-tune derived from it (Q8\_0). Each Qwen3.5 MoE layer also contains one always-active shared expert; the captured router logits cover only the 256 routed experts, and every expert set in this paper is over those. Qwen3.5-122B-A10B (HauhauCS, Q8\_K\_P) uses the same 256-expert, top-8 routing scheme with 48 MoE layers. GLM-4.7-Flash, official and HauhauCS, routes each token to 4 of 64 experts plus one shared expert by sigmoid gating with a per-layer correction bias.

For the Qwen checkpoints, router logits were captured at every MoE layer with a hook in llama.cpp under greedy decoding (temperature 0, top-k 1, fixed seed), storing routing only. Prompt and generated tokens were retained as separate blocks and were never combined. Generation blocks were cut at the first end-of-turn token when that token was available; the unrecoverable 122B exception is described under Provenance and scope. Generation was captured for both Qwen3.5-35B checkpoints. A 122B generation capture also exists but is treated as exploratory because its trimming failed. The GLM captures are prefill only and were taken from BF16 weights with a Transformers hook.

\paragraph{Expert-set construction.} At each Qwen MoE layer, the router scores all 256 experts, and the scores are converted to probabilities by a softmax. The 8 highest were kept and rescaled to sum to 1, matching the weights the model uses to mix those experts. From these we construct the routed weight $W$, selection rate $S$, and conditional weight $Q$. For each token set, $W$ was computed for every expert as the mean renormalized weight, with unselected experts contributing zero; $S$ is the fraction of tokens on which the expert was selected, $Q$ its mean weight when selected, and $W = S \cdot Q$. Layers were pooled by an unweighted mean over the per-layer $W$ vectors, indexed by expert number. Because expert index $e$ at one layer and expert index $e$ at another layer denote different modules, these pooled objects are layer-pooled expert-index profiles rather than single experts persisting across depth. The primary claims therefore concern recurrence in layer-pooled routing organization. Per-layer identification is not available for the packed designs, whose surviving per-token tensors are layer-averaged; on the independent probe, where per-layer profiles survive, the generation-time dispersion is present at every layer, with about 14 distinct subject winners per layer across all 40 layers (prefill 1.6), slightly stronger in the middle third and weakest at the first and last layers (Appendix~\ref{app:perlayer}). The layer-pooled profile therefore summarizes a signal that is flat across depth rather than concentrated in a few layers. Prompt and generation blocks were never pooled, because the prompt block would dominate any combined set. The expert-index set of a token set is its top 8 indices by layer-pooled $W$. Sets are compared by the Jaccard index, whose exact expectation for two random sets of 8 from 256 is 0.017 (the ratio approximation $k/(2E-k)$ gives 0.016). For GLM-4.7-Flash, the per-token rule is the model's own: the top 4 experts by sigmoid score plus the per-layer correction bias, weighted by their plain sigmoid values rescaled to sum to 1; the remaining steps are unchanged with 64 experts and sets of 4.

\paragraph{Prompt designs.} The independent probe asks about 20 subjects, with three questions per subject and one question per prompt; each answer was allowed up to 2,056 generated tokens. A subject's probe profile is computed over its three answers. The packed designs place all 20 questions in one prompt and let the model answer them in one generation. Three packed prompts were used: A (explanations), B (biographies), and C (synthesis), each padded to 446 prompt tokens and allowed up to 2,048 generated tokens. The three prompts list the subjects in the same order; in C the model regrouped its answer under thematic headers, so its segments occur in a different order from A. Answer spans were assigned to subjects at the model's own markdown headers, giving 12 segments in A, 15 in B, and 19 in C; the model skipped law in C and stopped early in A and B. Generation boundaries were taken from token IDs. No packed answer reached an end-of-turn token before the cap except C, which ended at token 1,812. A and C retain an expository register while subject changes. B changes the requested task and register while retaining the same subject labels, allowing register sensitivity to be measured across designs. For prefill spans, the Hugging Face tokenizer counts 391 to 392 tokens for the packed prompts where llama.cpp counted 446, so question spans were rescaled proportionally and may be off by one or two positions; generation spans are exact. The official instruction-tuned checkpoint design contains 12 prompts. Finance contributes two consequence-framed/neutral pairs (four prompts), while medical, physical, professional, and legal each contribute one pair (eight prompts); all answers were allowed 256 generated tokens. A 30-prompt set containing one long, coherent technical paragraph in three registers serves as a scoping analysis, not primary evidence.

\paragraph{Cross-design identification.} For each subject segment we form its top-8 set by layer-pooled $W$. A query segment in design X is compared by Jaccard with every candidate segment in design Y, and the predicted subject is the candidate with the highest similarity. Only subjects present in both designs enter either direction, and ties count as misses. Thus A$\leftrightarrow$P contributes $12 + 12 = 24$ queries and C$\leftrightarrow$P contributes $19 + 19 = 38$. The pooled generation analysis contains 168 bidirectional queries after intersecting the subjects available in each design; Table~\ref{tab:pairs} lists every ordered pair with its shared-subject count (A 12, B 15, C 19 and P 20 segments; law is absent from C, and A and B were capped by the generation budget). On prefill rows every design contains all 20 questions, so each pair contributes 40 queries and all six design pairs contribute 240. Identification accuracy is the fraction of queries whose predicted subject is correct; it is distinct from exact set recurrence, which we report separately as matched-subject Jaccard.

\paragraph{Controls.} Shuffled-label null: subject labels were permuted across candidates 200 times (NumPy \texttt{default\_rng(11)}) for generation and 100 times (seed 7) for prefill, and identification was recomputed; we report the mean, 95th percentile, and maximum. Position-matched null: A is compared with C, which the model answered in a different order, contrasting the same subject at different positions with different subjects at matched positions. Adjacent-segment null: candidates are restricted to the previous, same, and next subject in the target's order; first and last segments have two candidates rather than three and were not excluded, so 21 of the 168 queries have two candidates and the exact chance level is 0.354. For queries whose target is the independent probe, which has no answer order, the query design's order is used. Prefill control: the identical procedure is applied to the question-token rows of each design. Official-checkpoint replication: within-subject and between-subject Jaccard are compared across the crossed register conditions, followed by nearest-neighbour identification by subject among the other 11 prompts; because finance contributes four prompts, the chance level is 0.15. Set-size sensitivity: the prefill-versus-generation overlap was recomputed at top-8, top-15, and top-30; all three definitions gave the same qualitative conclusion.

\paragraph{Statistical analysis.} Jaccard similarity is $|A\cap B|/|A\cup B|$. Identification accuracies are counts over the denominators above and are reported with two-sided 95\% Wilson confidence intervals and against the shuffled-label distribution. Bidirectional queries share segments and are not independent; the intervals treat them as independent, and one-way counts are given in Table~\ref{tab:pairs}. No multiple-comparison correction is applied; the wording, position, drift, register, and shuffle controls were designed after the primary result and every one of them is reported. Within each captured segment, tokens contribute equally to $W$; layers contribute equally to the layer-pooled profile; and each eligible segment contributes one identification query regardless of its length. For the independent probe, $W$ was computed per answer and the subject profile is the unweighted mean over its three answers and 40 layers, so each answer carries equal weight regardless of its length.

\paragraph{Provenance and scope.} Checkpoint files, hashes where recorded, llama.cpp builds, capture commands, prompt hashes, and per-run metadata are listed in the data repository's \texttt{PROVENANCE.md}; the curated analysis repository is frozen at commit \texttt{f06cd3e}. The 35B fine-tune is SHA-256 \texttt{f3235db7\ldots}, its 60-prompt domain probe used llama.cpp build 8493 (\texttt{1772701f}), the 122B capture binary has SHA-256 \texttt{c3e205b3\ldots}, and the GLM register generation runs use commit \texttt{6658925}; the GLM domain prefill captures were taken with a Transformers hook and have no llama.cpp build. Per-token layer-averaged profiles survive for the packed prompts, and raw per-layer tensors survive for the controls and the GLM captures. The independent probe retains per-subject, per-layer summaries, and no 122B domain-probe tensors survive. Four exceptions determine the paper's scope. First, the 122B generation trimmer was a no-op, leaving approximately 45\% spill in its nominal generation block; that generation result is exploratory and is excluded from the primary replication claim. Second, the 122B generation-run prefill contains a last-layer artifact; a separate routing-only capture is clean and assigns 13 of 20 subjects to one expert. Third, GLM has no generation capture and supports only the prefill half of the analysis. Fourth, the official 35B checkpoint is instruction-tuned, so its comparison with HauhauCS tests persistence across that fine-tune, not emergence before instruction tuning. The primary packed-generation evidence therefore comes from the HauhauCS Qwen3.5-35B capture, with the smaller official instruction-tuned domain/register experiment serving as the generation-side checkpoint replication.

% CHANGELOG (Jeffrey's 2026-08-30 text -> this file). Everything not listed is verbatim.
% Abstract p3: "reordering them across prompts" -> "letting the model regroup them across prompts" (fixlist 1); added "prefill on all five and generation on the two Qwen3.5-35B-A3B checkpoints" (23).
% Abstract p4: added "when the answer's register is held constant (0.70 over all design pairs, 0.54 when the register changes)" (16); "to" for ranges.
% Intro p2, p3: [citation] -> \citep keys.
% Intro p4: "reorder the subjects across prompts" -> "use the model's own regrouping of that answer to place subjects at different positions across prompts" (1); en-dash -> "to".
% Intro p5: "Reordering places" -> "The model's regrouping places" (1); "supporting evidence about the shared prefill regime" -> "the prefill replication" (reframe decision).
% Intro p7: "against 33% by chance" -> "against 35% by chance" (2).
% Intro p9: contributions now four; new first line = prefill replication (12:59 decision); register sentence -> "second organizer ... lowering identification from 0.85 to 0.54" (22).
% Methods models: shared-expert sentence added (29). Paragraph headings as \paragraph.
% Methods expert-set: {AUTHOR} hole -> per-layer statement with appendix ref (A.1).
% Methods prompt designs: "The subject order in C differs from A." -> the model regrouped (1); "392" -> "391 to 392" (13).
% Methods cross-design: {AUTHOR} hole -> Table 2 pointer with segment counts (A.2).
% Methods controls: "whose subjects occur in a different order" -> "which the model answered in a different order" (1); {AUTHOR} hole -> edge handling, chance 0.354, probe-target order rule (A.3, 21); chance 0.15 for the 12-prompt design (5).
% Methods statistical: non-independence sentence added (18); "prespecified" -> "designed after the primary result and every one of them is reported" (11); {AUTHOR} hole -> equal-weight averaging (A.4).
% Methods provenance: build 8493 -> the fine-tune's probe; commit 6658925 -> GLM register generation runs, prefill has no build (3); "Per-token tensors survive for the packed prompts" -> layer-averaged profiles (4); "no 122B tensors" -> "no 122B domain-probe tensors" (4).
